\documentclass[11pt]{article}
\usepackage{amsmath,amssymb,amsthm,hyperref}
\newtheorem{theorem}{Theorem}
\newtheorem{lemma}{Lemma}
\newtheorem{corollary}{Corollary}
\newcommand{\N}{\mathbb N}
\newcommand{\ellog}{\ell}
\title{A conditional logarithmic barrier for Erd\H{o}s Problem 243}
\author{Research attempt prepared in the accompanying Codex conversation}
\date{6 November 2026}
\begin{document}
\maketitle

\begin{abstract}
This note proves a conditional rigidity criterion for the integer remainder
states associated with Erd\H{o}s Problem 243. If the positive increments of
the numerator are eventually at most
$\alpha\log\log C_n+O(1)$, with $\alpha\le1/\log 2$, the sequence must
eventually obey the Sylvester recurrence. The argument combines a bound
on the common divisor, control of prime reuse, and a Chinese remainder
barrier. A further sieve argument constrains the denominator prime support
of any counterexample with finite logarithmic rise coefficient.
The barrier also holds at the endpoint for the smaller LCM-cleared
numerator, even when redundant denominator factors grow without bound.
A clustered family of CRT barriers gives a further denominator-support
bound for its finite rise coefficient.
The unrestricted problem does not supply this increment bound.
This is not a proof or disproof of the full conjecture. Novelty of the
conditional criterion has not been established.
\end{abstract}

\section{Problem and exact states}
The \href{https://conjectures.io/problems/erdos243-erdos-243}{listed problem}
asks whether strictly increasing positive integers satisfying
\[
 a_{n+1}/a_n^2\longrightarrow1,
 \qquad S=\sum_{n\ge0}\frac1{a_n}\in\mathbb Q
\]
must satisfy $a_{n+1}=a_n^2-a_n+1$ eventually. The problem is also
\href{https://www.erdosproblems.com/243}{recorded as open by Erd\H{o}s Problems}.

Write $S=p/s$ with $s>0$, and put
\[
 R_n=\sum_{j\ge n}\frac1{a_j},\qquad
 D_n=s\prod_{j<n}a_j,\qquad C_n=D_nR_n.
\]
Then $C_n,D_n$ are positive integers, since
\[
 C_n=p\prod_{j<n}a_j-
 s\sum_{i<n}\prod_{\substack{j<n\\j\ne i}}a_j.
\]
They obey the exact recurrences
\begin{equation}\label{eq:state}
 D_{n+1}=a_nD_n,\qquad C_{n+1}=a_nC_n-D_n.
\end{equation}
These rational-tail coordinates and the associated centered residues are
developed in Koizumi's
\href{https://math.colgate.edu/~integers/aa28/aa28.pdf}{2026 paper}.
The earlier bounded-error CRT approach is described in Bado's
\href{https://www.researchgate.net/publication/413940355}{preprint}
and in Cook's
\href{https://github.com/wcook04/plectis-erdos/blob/main/paper/243/erdos-243-reciprocal-tail-rigidity.tex}{research note}.
The proof below allows a specified unbounded increment envelope.

The squaring hypothesis gives, after deleting finitely many terms,
$a_{n+1}\ge a_n^2/2\ge2a_n$. Consequently
\[
 R_n=\frac1{a_n}+O\!\left(\frac1{a_{n+1}}\right)
     =\frac{1+o(1)}{a_n}.
\]
It follows that
\begin{equation}\label{eq:ratio}
 \frac{C_{n+1}}{C_n}
 =\frac{a_nR_{n+1}}{R_n}
 =\frac{a_n^2}{a_{n+1}}(1+o(1))\longrightarrow1.
\end{equation}
Define, using natural logarithms,
\[
 \ellog(x)=\log\log\max\{x,\exp(e)\}.
\]
This only avoids irrelevant small-argument issues.

\section{Conditional criterion}
\begin{theorem}\label{thm:main}
Let $a_n\ge2$, $C_n\ge1$, and $D_n\ge1$ be integers satisfying
\eqref{eq:state}. Suppose
\[
 a_n\to\infty,\qquad a_{n+1}/a_n^2\to1,
 \qquad C_{n+1}/C_n\to1.
\]
If constants $K$ and $0\le\alpha<1/\log2$ exist such that eventually
\begin{equation}\label{eq:envelope}
 C_{n+1}-C_n\le\alpha\ellog(C_n)+K,
\end{equation}
then $C_n$ is eventually constant and
$a_{n+1}=a_n^2-a_n+1$ eventually.
\end{theorem}

\begin{proof}
Discard a finite prefix and increase $K$ so that $K\ge1$ and
\eqref{eq:envelope} holds at every remaining index. Reindex that tail.
Let $h(x)=\alpha\ellog(x)+K$, which is nondecreasing and positive.

First, a positive integer sequence with consecutive ratio tending to one
either tends to infinity or is eventually constant. Indeed, if it visits
$[1,M]$ at arbitrarily late indices, choose a time after which every ratio
differs from one by less than $1/(2M)$. A visit to $C_n\le M$ forces
$|C_{n+1}-C_n|<1/2$, hence equality. Induction traps it forever.

If $C_n=C$ eventually, \eqref{eq:state} gives $D_n=(a_n-1)C$ and
\[
 (a_{n+1}-1)C=D_{n+1}=a_nD_n=a_n(a_n-1)C,
\]
which proves the conclusion. Hence assume for contradiction that
$C_n\to\infty$.

\paragraph{The common divisor is bounded by the rise envelope.}
Put $G_n=\gcd(C_n,D_n)$. Equation \eqref{eq:state} implies
$G_n\mid G_{n+1}$. We claim
\begin{equation}\label{eq:gcd}
 G_n\le h(C_n)\quad\hbox{for every }n.
\end{equation}
If $G_n>h(C_n)$, take the first subsequent strict rise
$C_{t+1}>C_t$, $t\ge n$, which exists because $C_j\to\infty$.
Before this rise $C_t\le C_n$. Also $G_n$ divides
$C_{t+1}-C_t$, so
\[
 G_n\le C_{t+1}-C_t\le h(C_t)\le h(C_n)<G_n,
\]
a contradiction. Since $C_{n+1}/C_n\to1$, telescoping its logarithm
gives $\log C_n=o(n)$. In particular \eqref{eq:gcd} yields
\begin{equation}\label{eq:gcdsize}
 G_n=O(\log(n+2)).
\end{equation}

\paragraph{Almost every multiplier introduces entirely new prime support.}
Let
\[
 L_n=\operatorname{lcm}(D_0,a_0,\ldots,a_{n-1}),
 \qquad \Delta_n=D_n/L_n.
\]
The quotient is an integer, and
\[
 \frac{C_n}{D_n}=\frac{C_0}{D_0}-\sum_{j<n}\frac1{a_j}
\]
shows that $L_nC_n/D_n$ is an integer. Thus $\Delta_n\mid C_n$;
since it also divides $D_n$, we have $\Delta_n\mid G_n$.
Furthermore,
\begin{equation}\label{eq:reuse}
 \frac{\Delta_{n+1}}{\Delta_n}
 =\frac{a_nL_n}{\operatorname{lcm}(a_n,L_n)}
 =\gcd(a_n,L_n).
\end{equation}
Each index with $\gcd(a_n,L_n)>1$ multiplies $\Delta_n$ by at least two.
By \eqref{eq:gcdsize}, there are only $O(\log\log(N+3))$ such indices
below $N$. Enumerate all remaining indices as $i_0<i_1<\cdots$, and set
$b_j=a_{i_j}$. Then the $b_j$ are pairwise coprime and
\begin{equation}\label{eq:indices}
 i_j=j+O(\log\log(j+3)).
\end{equation}
The latter follows by subtracting the number of omitted indices from $i_j$;
first obtain $i_j=O(j+1)$ and then substitute this into the error bound.

\paragraph{The sizes of available moduli.}
Put $x_n=\log a_n$. The squaring hypothesis gives
$x_{n+1}=2x_n+\eta_n$, where $\eta_n\to0$.
For some $m,M>0$, therefore,
\begin{equation}\label{eq:doubleexp}
 m2^n\le\log a_n\le M2^n
\end{equation}
eventually. To justify positivity of $m$, start at $N$ so late that
$|\eta_n|\le1$ for $n\ge N$ and $x_N>1$.
Iteration gives $x_{N+k}\ge2^k(x_N-1)+1$.
The upper estimate follows in the same way.

For each large positive integer $B$, choose
$r_B=O(\log\log(B+3))$ such that $b_j>2B$ for every $j\ge r_B$.
This follows from the lower bound in \eqref{eq:doubleexp} and $i_j\ge j$;
finite exceptional indices are harmless. Choose the $B$ moduli
\[
 b_{r_B},b_{r_B+1},\ldots,b_{r_B+B-1}.
\]
Their largest original index is
$J_B=i_{r_B+B-1}=B+o(B)$ by \eqref{eq:indices}.
If $Q_B$ denotes their product, the upper bound in \eqref{eq:doubleexp}
and the distinctness of their indices give
\[
 \log Q_B\le M\sum_{i\le J_B}2^i\le M2^{J_B+1},
 \qquad
 \log\log Q_B\le(B+o(B))\log2.
\]

\paragraph{A forbidden block at a controlled height.}
Let $H_B=\max_{0\le n\le J_B+1}C_n$.
The relation $\log C_n=o(n)$ implies $\log H_B=o(B)$: for each
$\varepsilon>0$, all but finitely many terms satisfy
$\log C_n\le\varepsilon n$, and the finite maximum contributes $O(1)$.
The Chinese remainder theorem supplies an integer $T_B$ such that
\[
 T_B+j\equiv0\pmod{b_{r_B+j}}\quad(0\le j<B),
 \qquad H_B<T_B\le H_B+Q_B.
\]
Consequently
\[
 \log\log(T_B+B)\le(B+o(B))\log2.
\]
For example, this follows by bounding $T_B+B$ by $H_B+Q_B+B$ and
taking the maximum of the three corresponding logarithmic size bounds.
As $T_B>H_B\to\infty$ and $\alpha\log2<1$, it follows that
\begin{equation}\label{eq:width}
 h(T_B+B)<B
\end{equation}
for all sufficiently large $B$.

Let $N$ be the first index with $C_N\ge T_B$.
It exists because $C_n\to\infty$, and $N>J_B+1$ by the definition
of $H_B$. Since $C_{N-1}\le T_B-1$, we obtain
\[
 0\le C_N-T_B<h(T_B)<B.
\]
Thus $C_N=T_B+j$ for some integer $0\le j<B$.
The modulus $b=b_{r_B+j}$ divides both $C_N$ (by construction) and
$D_N$ (its original index is less than $N$). Using \eqref{eq:gcd}
and \eqref{eq:width},
\[
 2B<b\le G_N\le h(C_N)\le h(T_B+B)<B,
\]
the required contradiction.
\end{proof}

\begin{corollary}\label{cor:necessary}
For any counterexample to Erd\H{o}s Problem 243, with the integer states
defined above,
\[
 C_n\to\infty,\qquad
 \limsup_{n\to\infty}
 \frac{C_{n+1}-C_n}{\ellog(C_n)}\ge\frac1{\log2}.
\]
\end{corollary}
\begin{proof}
Bounded visits would force a Sylvester tail by the first paragraph of the
proof. If the displayed upper limit were smaller than $1/\log2$, choose
$\alpha$ strictly between it and $1/\log2$, with $\alpha\ge0$.
Theorem~\ref{thm:main} would apply, a contradiction.
\end{proof}

\section{A barrier for the LCM-cleared numerator}
The product-cleared numerator can include redundant denominator factors.
Removing those factors gives another conditional criterion, without a
bounded common-divisor hypothesis. The coordinates and their exact updates
are the LCM coordinates used by Bado and Cook, cited above; the bounded-rise
case is included in their earlier criteria.

\begin{theorem}\label{thm:lcmbarrier}
Assume the integer recurrences and limits of Theorem~\ref{thm:main}.
For every $n\ge0$ define
\[
 L_n=\operatorname{lcm}(D_0,a_0,\ldots,a_{n-1}),\quad
 \Delta_n=D_n/L_n,\quad U_n=C_n/\Delta_n.
\]
These $U_n$ are positive integers. If constants
$0\le\alpha<1/\log2$ and $K$ satisfy, eventually,
\[
 U_{n+1}-U_n\le\alpha\ellog(U_n)+K,
\]
then $C_n$ is eventually constant and the multipliers eventually satisfy
$a_{n+1}=a_n^2-a_n+1$.
\end{theorem}

\begin{proof}
The exact rational identity already used in the proof of
Theorem~\ref{thm:main} gives $\Delta_n\mid C_n$.
Put
\[
 E_n=D_n-(a_n-1)C_n,\quad
 V_n=E_n/\Delta_n=L_n-(a_n-1)U_n,\quad
 \rho_n=\gcd(a_n,L_n).
\]
Thus $V_n$ is an integer, and the exact updates are
\begin{equation}\label{eq:lcmupdates}
 \Delta_{n+1}=\rho_n\Delta_n,\qquad
 \rho_n U_{n+1}=U_n-V_n,\qquad
 \frac{V_n}{U_n}=1-\frac{C_{n+1}}{C_n}\longrightarrow0.
\end{equation}
In particular, once $|V_n|<U_n$, an index with $\rho_n\ge2$ is a strict
decrease of $U_n$. Every sufficiently late strict rise therefore has
$\rho_n=1$ and jump $U_{n+1}-U_n=-V_n$.

Suppose $C_n$ is not eventually constant. A late zero of $E_n$ would
give $C_{n+1}=C_n$ and $D_{n+1}/C_{n+1}\in\mathbb N$. At the next
step $C_{n+2}/C_{n+1}$ would then be a positive integer. Taking the
time sufficiently late that all subsequent consecutive ratios lie in
$(1/2,3/2)$ forces that integer to be one. Induction gives a constant
tail, a contradiction. Hence $E_n$, and therefore $V_n$, is nonzero
at every sufficiently late index. Since $V_n$ is integral and
$V_n/U_n\to0$, it follows that $U_n\to\infty$.

As before, $\log C_n=o(n)$. Since
$1\le\Delta_n\le C_n$, we have $\log\Delta_n=o(n)$.
The first update in \eqref{eq:lcmupdates} shows that the number of
indices $j<n$ with $\rho_j>1$ is $o(n)$: every such index multiplies
$\Delta$ by at least two. Enumerate the remaining indices as
$i_0<i_1<\cdots$, with $b_j=a_{i_j}$. Then
\[
 i_j=j+o(j),
\]
and the $b_j$ are pairwise coprime. We use the original $L_n$ throughout;
no change of its initial denominator is made.

For a large positive integer $B$, choose an offset
$r_B=O(\log\log(B+3))$ for which $b_j>2B$ whenever $j\ge r_B$.
The bounds $m2^n\le\log a_n\le M2^n$ established above justify this
choice. Select $b_{r_B},\ldots,b_{r_B+B-1}$, and put
\[
 J=i_{r_B+B-1}=B+o(B),\qquad
 Q=\prod_{j<B}b_{r_B+j},\qquad H=\max_{0\le n\le J+1}U_n.
\]
Since $U_n\le C_n$,
\[
 \log\log Q\le(B+o(B))\log2,\qquad \log H=o(B).
\]
Choose by CRT an integer $H<T\le H+Q$ such that
$b_{r_B+j}\mid T+j$ for $0\le j<B$, and set $\tau=T+B$.
For $h(x)=\alpha\ellog(x)+K$, increasing $K$ to be positive if needed,
these estimates give
\[
 h(\tau)<B
\]
for sufficiently large $B$.

Let $N$ be the first index with $U_N\ge\tau$. Then $N>J+1$ and
$U_{N-1}<\tau$. This is a sufficiently late rise, so
$\rho_{N-1}=1$ and its positive integer jump $d=U_N-U_{N-1}$ satisfies
\[
 0<d\le h(U_{N-1})\le h(\tau)<B.
\]
The starting point of the rise lies in $[T,T+B)$: indeed
$U_{N-1}=U_N-d\ge\tau-d>T$. Write $U_{N-1}=T+j$ with $0\le j<B$.
The corresponding $b=b_{r_B+j}$ divides both $U_{N-1}$ and $L_{N-1}$,
because its original index is at most $J<N-1$.
The identity $V_{N-1}=L_{N-1}-(a_{N-1}-1)U_{N-1}$ shows that $b\mid V_{N-1}$.
But $d=-V_{N-1}$, so $d\ge b>2B$, a contradiction.
Thus $C_n$ is eventually constant, and the asserted multiplier recurrence
follows as in Theorem~\ref{thm:main}.
\end{proof}

This condition is expressed in $U_n$, not in $C_n$. Positive changes of
$C_n=\Delta_nU_n$ can be much larger when $\Delta_n$ grows. The theorem
therefore does not assume the increment envelope of Theorem~\ref{thm:main}.
The unrestricted problem also does not currently supply the new envelope.

\begin{corollary}
For a counterexample to Erd\H{o}s Problem 243, the LCM-cleared numerator
$U_n=L_nR_n$ satisfies
\[
 U_n\longrightarrow\infty,\qquad
 \limsup_{n\to\infty}\frac{U_{n+1}-U_n}{\ellog(U_n)}\ge\frac1{\log2}.
\]
\end{corollary}
\begin{proof}
The divergence was proved above. An upper limit smaller than $1/\log2$
would supply the envelope in Theorem~\ref{thm:lcmbarrier}.
\end{proof}

\section{Finite rise coefficient and a sieve refinement}
The strict inequality in Theorem~\ref{thm:main} is enough to obtain a
stronger necessary condition, including the endpoint. The coordinates
in the next statement are fixed; rescaling them rescales the coefficient.

\begin{theorem}\label{thm:wheel}
Assume the integer recurrences and limits of Theorem~\ref{thm:main},
and suppose that $C_n$ is not eventually constant. Define
\[
 \Lambda=\limsup_{n\to\infty}
 \frac{C_{n+1}-C_n}{\ellog(C_n)}.
\]
If $\Lambda<\infty$, then $G_n=\gcd(C_n,D_n)$ is eventually a fixed
integer $g\ge1$. On that tail put $c_n=C_n/g$ and $d_n=D_n/g$, and let
\[
 \mathcal P=\{p\text{ prime}:p\mid d_n\text{ for some sufficiently late }n\}.
\]
For every finite $P\subseteq\mathcal P$, write
$\rho(P)=\prod_{p\in P}(1-1/p)$, with the empty product equal to one.
Then
\begin{equation}\label{eq:wheelbound}
 \Lambda\ge\frac{g}{\rho(P)\log2}.
\end{equation}
In particular
\[
 \prod_{p\in\mathcal P}(1-1/p)>0,\qquad
 \sum_{p\in\mathcal P}\frac1p<\infty,
 \qquad \Lambda>\frac1{\log2}.
\]
\end{theorem}

\begin{proof}
The integer-ratio argument already proved shows that $C_n\to\infty$.
Corollary~\ref{cor:necessary} gives $\Lambda\ge1/\log2$.
For each fixed index $N$, the integer $G_N$ divides every later $C_n,D_n$.
Dividing both coordinates by $G_N$ preserves all the recurrences and
limits, and $\ellog(C_n/G_N)/\ellog(C_n)\to1$. The new upper limit
of normalized rises is therefore $\Lambda/G_N$. If
$G_N>\Lambda\log2$, Theorem~\ref{thm:main} applied to this divided tail
would force eventual constancy, a contradiction. Thus $G_N\le\Lambda\log2$
for every $N$. The divisibility chain $G_N\mid G_{N+1}$ is bounded and
must eventually be constant, say equal to $g$.

The normalized coordinates satisfy $\gcd(c_n,d_n)=1$ at every later
index. If a prime divided both $a_n$ and $d_n$, it would divide both
$c_{n+1}=a_nc_n-d_n$ and $d_{n+1}=a_nd_n$, a contradiction.
Hence $\gcd(a_n,d_n)=1$. In particular the late $a_n$ are pairwise
coprime and are coprime to the denominator at the start of any such tail.
Write $\lambda=\Lambda/g$.

Fix finite $P\subseteq\mathcal P$, put $M=\prod_{p\in P}p$ and
$\rho=\rho(P)$. Reindex a fixed sufficiently late tail so that
$M\mid d_0$ and the coordinates are primitive throughout.
Suppose, towards a contradiction, that $\lambda\rho\log2<1$.
Choose $\beta>\lambda$ with $\beta\rho\log2<1$, and $K\ge1$ so that
\[
 c_{n+1}-c_n\le h(c_n),\qquad h(x)=\beta\ellog(x)+K
\]
on that tail (increasing $K$ handles its finite prefix).

Let $W$ tend to infinity through positive multiples of $M$.
There are exactly $B=\rho W$ integers $0\le j<W$ coprime to $M$;
list them as $j_0,\ldots,j_{B-1}$. When $M=1$, every $j$ is included.
Use the $B$ moduli $a_0,\ldots,a_{B-1}$, which are pairwise coprime
and coprime to $M$. Put
\[
 Q=M\prod_{i<B}a_i,\qquad H=\max_{0\le n\le B}c_n.
\]
The growth estimates proved above give
\[
 \log\log Q\le B\log2+O(1),\qquad \log H=o(B).
\]
The Chinese remainder theorem supplies $H<T\le H+Q$ satisfying
\[
 T\equiv0\pmod M,\qquad T+j_i\equiv0\pmod{a_i}\quad(i<B).
\]
Consequently $\ellog(T+W)\le B\log2+O(1)$ and
\[
 h(T+W)\le\beta\rho W\log2+O(1)<W
\]
for sufficiently large $W$.

At the first index $N$ with $c_N\ge T$ we have $N>B$ and
$0\le c_N-T<h(T)<W$. Write $c_N=T+j$ with $0\le j<W$.
If $\gcd(j,M)>1$, a prime of $M$ divides $c_N$ and $d_N$.
Otherwise $j=j_i$ for some $i<B$, and $a_i$ divides both $c_N$ and $d_N$.
Either case contradicts $\gcd(c_N,d_N)=1$. This proves
$\lambda\rho\log2\ge1$, equivalently \eqref{eq:wheelbound}.

Since \eqref{eq:wheelbound} holds for every finite $P$, every such
product satisfies $\rho(P)\ge g/(\Lambda\log2)>0$. Thus the product
over $\mathcal P$ is positive. The inequality
$-\log(1-1/p)\ge1/p$ then proves convergence of the reciprocal sum.
Finally $\mathcal P$ is nonempty: $d_{n+1}=a_nd_n$ and $a_n\ge2$.
Using a singleton $P=\{p\}$ in \eqref{eq:wheelbound} gives
$\Lambda\ge g/((1-1/p)\log2)>1/\log2$.
\end{proof}

\begin{corollary}\label{cor:endpoint}
The conclusion of Theorem~\ref{thm:main} also holds when
$\alpha=1/\log2$.
\end{corollary}
\begin{proof}
If $C_n$ were not eventually constant, its rise bound would give
$\Lambda\le1/\log2$, contrary to Theorem~\ref{thm:wheel}.
\end{proof}

\section{Clustered barriers for the LCM-cleared numerator}
The following refinement does not require the common divisor of $C_n,D_n$
to stabilize. We keep the original LCM coordinates of
Theorem~\ref{thm:lcmbarrier} throughout.

\begin{theorem}\label{thm:lcmwheel}
Assume the integer recurrences and limits of Theorem~\ref{thm:main},
and suppose that $C_n$ is not eventually constant. Let
\[
 \Lambda_U=\limsup_{n\to\infty}
 \frac{U_{n+1}-U_n}{\ellog(U_n)},\qquad
 \mathcal P_L=\{p\text{ prime}:p\mid L_n\text{ for some }n\}.
\]
For every finite $P\subseteq\mathcal P_L$, with
$\rho(P)=\prod_{p\in P}(1-1/p)$, one has
\[
 \Lambda_U\ge\frac1{\rho(P)\log2}.
\]
Consequently $\Lambda_U>1/\log2$. If $\Lambda_U<\infty$, then
\[
 \prod_{p\in\mathcal P_L}(1-1/p)>0,
 \qquad \sum_{p\in\mathcal P_L}\frac1p<\infty.
\]
In particular, if eventually
\[
 U_{n+1}-U_n\le\frac{\ellog(U_n)}{\log2}+K
\]
for a fixed $K$, then $C_n$ is eventually constant and
$a_{n+1}=a_n^2-a_n+1$ eventually.
\end{theorem}

\begin{proof}
We use the previously proved facts $U_n\to\infty$, $V_n/U_n\to0$,
$\log C_n=o(n)$, and
\[
 \rho_n U_{n+1}=U_n-V_n,\qquad
 L_{n+1}=a_nL_n/\rho_n,\qquad \rho_n=\gcd(a_n,L_n).
\]
Also $U_{n+1}/U_n=(C_{n+1}/C_n)/\rho_n$. Hence eventually
\begin{equation}\label{eq:reusecontraction}
 \rho_n>1\ \Longrightarrow\ U_{n+1}\le\tfrac34 U_n.
\end{equation}
If $\rho_n=1$ and $M\mid L_n$, then
\begin{equation}\label{eq:gcdclass}
 \gcd(U_{n+1},M)=\gcd(a_nU_n-L_n,M)=\gcd(U_n,M).
\end{equation}
Thus the gcd class is preserved between reuse steps.

Fix finite $P$, put $M=\prod_{p\in P}p$ and $\rho=\rho(P)$.
Suppose for contradiction that $\Lambda_U\rho\log2<1$.
Choose $\beta>\max\{\Lambda_U,0\}$ with $\beta\rho\log2<1$.
After discarding a fixed prefix, we may assume $M\mid L_0$,
\eqref{eq:reusecontraction}, and
\[
 U_{n+1}-U_n\le h(U_n),\qquad h(x)=\beta\ellog(x)+K,
 \quad K\ge1.
\]
The discarded prefix does not change the definitions of the LCM states.

Let $W$ tend to infinity through multiples of $M$ and put
$B=\rho W$. For every divisor $d\mid M$ define
\[
 J_d=\{0\le j<W:\gcd(j,M)=d\}.
\]
Then $|J_d|=W\varphi(M/d)/M\le B$, since $M$ is squarefree.
For $M=1$ there is just one class, containing every position.
Enumerate the fresh multipliers as in Theorem~\ref{thm:lcmbarrier}.
Select $B$ of them, denoted $b_1,\ldots,b_B$, all greater than $2W$.
Their last original index is $J=B+o(B)$, they are pairwise coprime,
and they are coprime to $M$. Set
\[
 Q=M\prod_{i=1}^B b_i,\qquad S=Q^2,\qquad
 H=\max_{0\le n\le J+1}U_n.
\]
The same growth estimates give
\[
 \log\log Q\le(B+o(B))\log2,
 \qquad \log H=o(B).
\]
In particular $S>H$ for large $W$, since $Q>(2W)^B$.

For each $d\mid M$, assign the positions of $J_d$ injectively to
the same list $b_1,\ldots,b_B$. The Chinese remainder theorem supplies
$0\le r_d<Q$ such that
\[
 r_d\equiv0\pmod M,\qquad
 r_d+j\equiv0\pmod{b_i}
 \quad\text{whenever position }j\in J_d\text{ is assigned to }b_i.
\]
Unused moduli may be given arbitrary residues. Put $T_d=S+r_d$.
These barriers all lie in the common interval
$[S,S+Q+W]$. Their construction uses $B$ shared moduli,
rather than a separate set of moduli for every class. Moreover
\[
 h(S+Q+W)\le\beta(B+o(B))\log2+O(1)<W.
\]

Let $N$ be the first index with $U_N\ge S+Q+W$, and let $k<N$
be the last index with $U_k<S$. Both exist because $U_n\to\infty$
and $U_0<S$. We have $k\ge J+1$ and
\[
 S\le U_{k+1}<S+h(S)<S+W.
\]
For $k+1\le n<N$, a reuse step would give
\[
 U_{n+1}\le\tfrac34 U_n
 <\tfrac34(S+Q+W)<S,
\]
contradicting the choice of $k$. Thus $\rho_n=1$ throughout this
crossing and $\gcd(U_n,M)$ has a fixed value $d$ there.

We start below $T_d+W$ and finish above it. At the first rise crossing
$T_d+W$, its starting point has the form $U_n=T_d+j$, $0\le j<W$,
and its positive jump $t$ is less than $W$. Since $T_d\equiv0\pmod M$,
we have $j\in J_d$. Its assigned modulus $b_i$ divides both $U_n$
and $L_n$, because $n>J$. As $\rho_n=1$,
\[
 t=U_{n+1}-U_n=(a_n-1)U_n-L_n,
\]
so $b_i\mid t$. This contradicts $0<t<W<b_i$.
The claimed coefficient bound follows.

The set $\mathcal P_L$ is nonempty since each $a_n\ge2$ divides
$L_{n+1}$. A singleton $P$ therefore gives $\Lambda_U>1/\log2$.
If $\Lambda_U$ is finite, the bound for all finite $P$ implies
$\rho(P)\ge1/(\Lambda_U\log2)>0$, which proves the Euler-product
and reciprocal-sum claims. Finally the endpoint envelope would imply
$\Lambda_U\le1/\log2$, a contradiction.
\end{proof}

\section{A failed Egyptian-length descent}
For a positive rational $r$, let $\lambda(r)$ be the least number of
positive unit fractions summing to $r$; repetitions are allowed.
It is tempting to use $\lambda$ as a decreasing integer invariant for
the near-Sylvester step $r\mapsto r-1/a$.
The following exact examples exclude a descent assertion based only
on the displayed local data. They do not exclude an argument using
additional information about all later states.

\begin{lemma}\label{lem:lengthfailure}
There is a positive rational $r=C/D$ and an integer $a\ge2$ with
\[
 \lambda(r)=2,\quad \lambda(r-1/a)\ge3,\quad
 0<\frac{aC-D-C}{C}<1.2\cdot10^{-14},
 \quad a=1+\operatorname{round}(D/C).
\]
There is a second pair with $\lambda(r)=2$, $\lambda(r-1/a)\ge3$,
$a=1+\operatorname{round}(D/C)$, $a>2\cdot10^7 C$,
$\gcd(a,D)=1$, and relative numerator rise less than $10^{-3}$.
\end{lemma}

\begin{proof}
For the first assertion take
\[
 b=29719973,\quad c=441638367840389,\quad a=29719972,
 \qquad C=b+c,\quad D=bc.
\]
The integers $a/4=7429993$, $b$, and $c$ are prime. They may be
checked by the exact trial divisions in the accompanying script.
Direct multiplication gives $aC-D=C+5$, so the relative rise is
$5/441638397560362<1.2\cdot10^{-14}$ and $a-1$ is the nearest integer
to $D/C$. Since $b,c$ are distinct primes, $r=1/b+1/c$ has length two.

Write $u=C+5$ and $v=aD$, so $r-1/a=u/v$, and put $M=5103$.
We have
\[
 M\mid u,\qquad v\equiv1030\pmod M,
 \quad a=4p,\quad p\equiv25,\ b\equiv101,\ c\equiv4997\pmod M.
\]
If $u/v=1/x+1/y$, then
\[
 (ux-v)(uy-v)=v^2,
 \qquad ux-v>0,\quad uy-v>0.
\]
Thus a positive divisor of $v^2$ must be congruent to $-v=4073$ modulo
$M$. Every divisor has the form $2^e p^f b^g c^h$ with
$0\le e\le4$, $0\le f,g,h\le2$. The following finite residue sets
are disjoint:
\begin{align*}
 \{2^e25^f\bmod M\}
 &=\{1,2,4,8,16,25,50,100,200,400,625,1250,2500,4897,5000\},\\
 \{4073(101^g4997^h)^{-1}\bmod M\}
 &=\{106,206,394,2020,3062,3082,4073,5003,5102\}.
\end{align*}
All inverses exist modulo $M$. This excludes two unit fractions,
and also one, since a unit fraction is the sum of two equal unit fractions.

For the second assertion put
\[
 L=1222709419,\quad a=121047033277,\qquad
 r=\frac1{101L}+\frac1{4997L}=\frac{5098}{504697L}.
\]
Both $L$ and $a$ are prime, and
\[
 L=1+5098\cdot239841,\qquad
 a=100+504697\cdot239841.
\]
Consequently $aC-D=C+5=5103$, $a>2\cdot10^7C$, and
$\gcd(a,D)=1$. The relative rise is $5/5098<10^{-3}$.
The input again has length two. The output denominator is
$v=19\cdot101\cdot263\cdot L\cdot a$, with $L\equiv1$ and
$a\equiv100$ modulo $M$; also $v\equiv1030\pmod M$.
The same divisor test now checks the $3^5=243$ positive divisors
of $v^2$. The script exhausts them and finds none congruent to $4073$
modulo $M$, proving the asserted length increase.
\end{proof}

\section{Verification and the unresolved step}
The script \texttt{erdos243-check.mjs} checks exact state and LCM identities
on all 3,044 reduced rationals $0<p\le q\le100$. It also constructs and
checks twelve exact CRT blocks, using pairwise coprime Sylvester numbers
as sample moduli, together with sieve blocks whose uncovered positions
are assigned separate coprime moduli. These are finite checks, not evidence of an infinite
termination theorem. Run the script with \texttt{node erdos243-check.mjs}.
The script also constructs clustered barriers for every gcd class
modulo $6$ and $30$, with one shared list of fresh moduli per cluster.
It verifies the prime factorizations, unit-fraction identities and all
divisor exclusions used in Lemma~\ref{lem:lengthfailure}.

An independent agent reviewed the conditional argument without finding a
substantive error. A separate Codex CLI review, using its configured
\texttt{gpt-6.1-sol} model, reached the same finding. Its report is saved as
\texttt{erdos243-codex-review.txt}; the supplied artifact and contract are
in \texttt{erdos243-review-prompt.txt}. Neither review is a formal proof
or a novelty certification. Those two reviews covered the strict-coefficient
theorem. A further fresh-context adversarial review found no substantive
error in the sieve refinement and endpoint corollary, assuming the earlier
strict-coefficient theorem. That report is in
\texttt{erdos243-sieve-review.txt}. A separately authorized Codex CLI
review also found no substantive error or invalidating unstated assumption
in the sieve theorem and endpoint argument, assuming the earlier strict
theorem; see \texttt{erdos243-sieve-codex-review.txt}. Its isolated input is
\texttt{erdos243-sieve-review-prompt.txt}.
The LCM-numerator theorem received a separate fresh-context adversarial
review, which found no substantive mathematical error under its stated
hypotheses. That report is in \texttt{erdos243-lcm-review.txt}.
No cross-model review of that theorem was run during the automatic goal
continuation.
The clustered LCM barrier in Theorem~\ref{thm:lcmwheel} received its own
fresh-context adversarial review under the isolated preliminary facts;
the reviewer found no substantive mathematical issue. Its report and
input are \texttt{erdos243-cluster-review.txt} and
\texttt{erdos243-cluster-review-prompt.txt}. This review does not
establish novelty or solve the unrestricted conjecture.
A separately authorized Codex CLI review found no substantive error
in the explicit clustered-barrier claims under the supplied preliminary
facts. It could not assess the endpoint corollary from its excerpt,
because the referenced earlier envelope was not included. The endpoint
envelope is now written explicitly in Theorem~\ref{thm:lcmwheel};
it implies $\Lambda_U\le1/\log2$ directly. The original isolated
review input is retained unchanged, and the review is in
\texttt{erdos243-cluster-codex-review.txt}.
The finite Egyptian-length counterexamples were separately reviewed
against their exact verifier; the reviewer found no substantive issue.
The isolated contract and report are
\texttt{erdos243-descent-review-prompt.txt} and
\texttt{erdos243-descent-review.txt}. Cross-model review was skipped
for this automatic goal continuation.

The original hypotheses yield only $C_{n+1}-C_n=o(C_n)$, and
\eqref{eq:envelope} is much stronger. No argument here closes that gap.
In particular, subexponential growth of an integer sequence does not
imply this increment bound. Excluding the remaining larger-increment
integer orbits is the work still required for the full problem.
For the LCM-cleared numerator, the known estimate on a positive jump is
only $U_{n+1}-U_n=o(U_n)$; this follows from \eqref{eq:lcmupdates}
and $\rho_n=1$ at late rises. The new logarithmic envelope has not been
deduced from that estimate or from the remaining arithmetic constraints.
\end{document}
